\begin{historical}
Roald Hoffmann was born Roald Safran in Złoczów, Poland, in 1937. When he was five, his father Hillel smuggled him and his mother out of a Nazi labour camp. A Ukrainian neighbour, Mykola Dyuk, hid them for eighteen months in the attic and storeroom of a local schoolhouse. His mother taught him to read using textbooks stored in the attic. His father remained in the camp and was killed in 1943 for organising a prisoner escape. Most of the extended family perished. Hoffmann and his mother survived, emigrated to the United States in 1949, and he eventually became a theoretical chemist at Cornell.

In 1952, Kenichi Fukui introduced frontier molecular orbital theory, identifying the HOMO and LUMO as the orbitals governing chemical reactivity. His framework was qualitative but pointed towards a deeper principle. Meanwhile, pericyclic reactions—concerted transformations such as electrocyclic ring closures and sigmatropic shifts—presented puzzling stereospecific behaviour that resisted classical explanation. Some worked thermally but not photochemically; others required light. The outcomes seemed predictable only in hindsight.

The breakthrough came from an anomaly. In the early 1960s, Robert Burns Woodward—already regarded as the greatest synthetic organic chemist of his generation, and a 1965 Nobel laureate—was pursuing the total synthesis of vitamin B\textsubscript{12}, a molecule with 181 atoms, 9 stereocenters, and a corrin ring that had resisted synthesis for over a decade. A critical electrocyclic step gave, in Woodward's words, \QSOpen{}precisely the opposite stereochemical course to that which we had predicted.\QSClose{} He had used physical CPK models, the standard tool of the era, to predict the stereochemistry and got it exactly backwards. Rather than dismissing the result, Woodward brought the puzzle to Hoffmann, who was then developing orbital phase methods using the extended Hückel approach. Their investigation of why the reaction went the \QSOpen{}wrong\QSClose{} way led directly to a series of papers (1965–69) articulating what became the Woodward–Hoffmann rules: pericyclic reactions follow strict constraints based on conservation of orbital symmetry. Whether a reaction is allowed or forbidden can be deduced by tracking how the symmetries of occupied orbitals evolve along the reaction coordinate.

Experimentalists immediately tested the rules across electrocyclic closures, sigmatropic shifts, and cycloadditions. Every known case conformed. The rules offered not just post hoc explanation but genuine predictive power, and by the early 1970s orbital symmetry had become a central organising principle in mechanistic organic chemistry. In 1981, Hoffmann shared the Nobel Prize in Chemistry with Fukui for their independent theoretical contributions to reaction mechanisms. Woodward had died in 1979, two years too early. Had he lived, he would have been among the very few scientists to win two Nobel Prizes.
\end{historical}
